\section{Energy of the central soft regional operator}

The following lemmas construct the central soft replacement, its regional
coordinates and a bound on its energy in the original ground space.
They are auxiliaries to the generation and energy-estimation steps of
\cite[Figures \texttt{fig:generate} and \texttt{fig:algo1}]{Arad2016RigorousRG}.
The harmonic cap retains the correction in
\cite{gap:alvv16_soft_truncation_log_bound}. The one-sided center and
scalar-centering distinctions are recorded in
\cite{gap:alvv16_soft_regional_energy_centering}.

\begin{definition}[Interactions restricted to an interval]
    \label{def:interval_bond_restriction}
    \lean{TNLean.MPS.AreaLaw.intervalBondRestriction}
    \leanok
    For a chain with $L$ sites, define $h_i^{[l,r)}=h_i$ if
    $l\le i$ and $i+1<r$, and $h_i^{[l,r)}=0$ otherwise.
\end{definition}
\begin{lemma}[The actual internal interaction sum]
    \label{lem:interval_bond_restriction}
    \lean{TNLean.MPS.AreaLaw.intervalBondRestriction_supported,
        TNLean.MPS.AreaLaw.intervalBondRestriction_bounds,
        TNLean.MPS.AreaLaw.sum_intervalBondRestriction,
        TNLean.MPS.AreaLaw.sum_internal_intervalBondRestriction_eq,
        TNLean.MPS.AreaLaw.sum_exterior_intervalBondRestriction}
    \leanok
    \uses{def:interval_bond_restriction}
    The restricted interactions retain their designated two-site supports
    and the bounds $0\le h_i^{[l,r)}\le I$. Their full sum is the original
    internal sum on $[l,r)$. Every subinterval of $[l,r)$ has the original
    internal sum; the exterior intervals $[0,l)$ and $[r,L)$ have zero sum.
\end{lemma}
\begin{proof}\leanok
    A retained bond is unchanged and every other bond is zero. The stated
    identities follow by separating the finite sums into these two cases.
\end{proof}

\begin{lemma}[The minimum and norm of the central replacement]
    \label{lem:soft_parent_structure}
    \lean{TNLean.MPS.AreaLaw.ALVV.exists_soft_parent_structure}
    \leanok
    \uses{lem:interval_bond_restriction,lem:alvv_five_soft_residuals}
    Suppose $0\le h_i\le I$, each $h_i$ is supported on sites $i,i+1$,
    and the local dimension is positive. Fix $l\le u\le v\le r\le L$
    with $u-l,r-v\le2\ell$. Let $P,M$ be the original sums internal
    to $[l,r),[u,v)$, respectively, and write $e_P,e_M$ for their minima.
    These attained minima are chosen before $p,t,m$.
    For $e_M-10\le p\le e_M$, $t>0$ and $m\ge2$, define
    \begin{align}
      R&=P-M+pI+h_{m,t}(M-pI),
      \label{eq:soft_parent_operator}\\
      B_*&=t(1+\log m)+4\ell+14.
      \label{eq:soft_parent_budget}
    \end{align}
    The actual attained minimum $e_{\mathrm s}$ satisfies
    $R-e_{\mathrm s}I\ge0$ and $\|R-e_{\mathrm s}I\|\le B_*$.
    It is chosen before every $T$. If $t,T\ge48$, then
    \begin{align}
      0\le e_P-e_{\mathrm s}&\le\rho_{m,t}(T),&
      \rho_{m,t}(T)&=32e^{-(T-48)/8}
             +\frac{t}{m}\left(\frac{T}{t}\right)^m.
      \label{eq:soft_parent_minimum}
    \end{align}
\end{lemma}
\begin{proof}\leanok
    Restrict the interactions to the parent interval and divide the
    chain at $0,l,u,v,r,L$. The exterior sums vanish. Choose attained
    minima of the five disjoint regional sums. Their sum differs from
    the parent minimum by a number in $[0,4]$. Apply the simultaneous
    regional soft residual estimate with only the middle region selected.
    Its center defect is at most ten and its crossing set has at most
    four elements. The reverse residual at its attained soft ground
    vector gives (\ref{eq:soft_parent_minimum}). The two collars contribute
    at most $4\ell I$, the central soft term contributes at most
    $t(1+\log m)I$, and the crossing terms contribute at most $4I$.
    The centered soft minimum is at least $-10$, giving
    (\ref{eq:soft_parent_budget}).
\end{proof}

\begin{lemma}[Regional coordinates and original-ground energy]
    \label{lem:native_soft_parent_energy}
    \lean{TNLean.MPS.AreaLaw.ALVV.exists_native_soft_parent_energy_bound}
    \leanok
    \uses{lem:soft_parent_structure,lem:regional_ground_energy_bound}
    Under the hypotheses of Lemma~\ref{lem:soft_parent_structure}, there
    is an actual positive matrix $K$ on the Hilbert space of $[l,r)$
    such that $R=K\otimes I+e_{\mathrm s}I$ and $\|K\|\le B_*$.
    These data are chosen before $T$ and every candidate space. For
    $t,T\ge48$ and every vector $g$ in the full ground space of the
    original $H=\sum_i h_i$,
    \begin{align}
       \langle g,(K\otimes I)g\rangle
          \le(2+\rho_{m,t}(T))\|g\|^2.
    \end{align}
    No gap or uniqueness hypothesis is imposed; empty intervals and
    chain endpoints are included.
\end{lemma}
\begin{proof}\leanok
    The centered replacement is positive and supported on $[l,r)$.
    Restriction to a fixed complementary configuration gives a positive
    regional matrix with the same upper norm bound. The downward soft
    order and (\ref{eq:soft_parent_minimum}) give
    $R-e_{\mathrm s}I\le P-e_PI+\rho_{m,t}(T)I$.
    On the original ground space the centered regional energy is at
    most two. Adding the scalar term proves the assertion.
\end{proof}

The energy bound concerns the original ground space. Applying it to the
separate untrimmed energy candidate in the generation procedure requires
an additional quantitative transfer to the same generated soft target.

\begin{lemma}[A middle center from a parent estimate]
    \label{lem:regional_central_energy_estimate}
    \lean{TNLean.MPS.AreaLaw.exists_regional_central_energy_estimate}
    \leanok
    \uses{lem:adjacent_interval_energy_merge}
    Let $0\le h_i\le I$ have their designated nearest-neighbour supports
    on a chain of positive local dimension, and fix
    $l\le u\le v\le r\le L$. The actual internal Hamiltonians on
    $[l,u)$, $[u,v)$, $[v,r)$ and $[l,r)$ have attained minima
    $e_{\mathrm L},e_{\mathrm M},e_{\mathrm R},e_{\mathrm P}$, respectively,
    with normalized attaining vectors and positive centered matrices.
    These data are chosen before every parent estimate, and satisfy
    \begin{align}
       e_{\mathrm L}+e_{\mathrm M}+e_{\mathrm R}
       \le e_{\mathrm P}
       \le e_{\mathrm L}+e_{\mathrm M}+e_{\mathrm R}+2.
    \end{align}
    If $e_{\mathrm P}-8\le p_{\mathrm P}\le e_{\mathrm P}$, then
    $p_{\mathrm M}=p_{\mathrm P}-e_{\mathrm L}-e_{\mathrm R}-2$
    satisfies $e_{\mathrm M}-10\le p_{\mathrm M}\le e_{\mathrm M}$.
    Empty intervals and endpoint cuts are included.
\end{lemma}
\begin{proof}\leanok
    Join $[l,u)$ to $[u,r)$, and then join $[u,v)$ to $[v,r)$.
    Each joining raises the sum of the two minima by a number in $[0,1]$.
    The two constructions of the minimum on $[u,r)$ agree, since both
    are attained positive-semidefinite lower bounds for the same matrix.
    Adding the two comparisons proves the bound with two. Substitution
    of the stated parent-estimate interval gives the middle-estimate
    interval.
\end{proof}

The middle-center conversion does not require a collar-length bound.
When the two collars have length at most $2\ell$, their original internal
matrices act on those short regional Hilbert spaces; their minima are
exactly those used here. No running-time conclusion for computing these
minima is asserted. This gives a corrected one-sided route to the soft
center, without identifying it with the source's broader absolute-error
input.
